% !TEX root = ../main.tex
\section{CSP Formulation and Implementation}

The formulation (variables, domains, three constraint types) is given in Section~1.
\code{Constraints} stores \code{variables} (file order), \code{num\_slots}, \code{precedence} pairs
and the \code{latest} map. Initial domains are $[1,\dots,T]$ for every variable. Before search
starts, \code{backtracking\_search} applies the unary latest-slot constraint to each domain (node
consistency). If any domain becomes empty, it returns \code{None} at once.

\textbf{The five required functions.}
\begin{itemize}
  \item \code{select\_unassigned\_variable} applies \textbf{MRV}: it picks the unassigned variable
    with the smallest current domain. Ties are broken by the \textbf{degree heuristic}: more
    \emph{distinct other unassigned variables} linked by a precedence constraint wins. Exclusivity
    links every pair equally, so it is not counted. Ties left after that go to file order
    D1\dots Dn (\dec{05}).
  \item \code{order\_domain\_values} applies \textbf{LCV} (optional in the spec, implemented). It
    sorts the variable's slots by how many values they would remove from the domains of other
    unassigned variables, fewest first. Python's stable sort keeps ascending slot order among ties.
  \item \code{is\_consistent} checks the proposed value against the latest-slot limit and against
    every \emph{already assigned} variable: no shared slot (exclusivity) and correct order for any
    precedence pair between them. Constraints with unassigned variables are left for forward
    checking.
  \item \code{forward\_check} copies the domains and, for each unassigned variable, removes every
    slot that would clash with the new assignment (same slot, or wrong side of a precedence pair).
    It returns the pruned copy, or \code{None} if some domain becomes empty (a dead end found early).
  \item \code{backtracking\_search} runs the recursive search below and returns a complete
    assignment, or \code{None} if no feasible schedule exists.
\end{itemize}

\textbf{Restoration.} \code{forward\_check} never changes the domains it is given; it works on a
copy. When a value fails, the recursion simply continues with the parent's unchanged domains. So
every value removed by forward checking is automatically ``restored'' on backtracking, and no undo
log is needed.

\begin{algorithm}[H]
\caption{Backtracking search with MRV, degree, LCV and forward checking}
\begin{algorithmic}[1]
\Function{Backtrack}{$assignment, domains, C$}
  \State $X \gets \Call{Select-Unassigned-Variable}{assignment, domains, C}$ \Comment{MRV, degree, file order}
  \If{$X = \textit{none}$} \Return copy of $assignment$ \Comment{all tasks scheduled}
  \EndIf
  \For{$v \in \Call{Order-Domain-Values}{X, assignment, domains, C}$} \Comment{LCV}
    \If{\textbf{not} \Call{Is-Consistent}{$X, v, assignment, C$}} \textbf{continue}
    \EndIf
    \State $assignment[X] \gets v$
    \State $pruned \gets \Call{Forward-Check}{X, v, assignment, domains, C}$ \Comment{works on a copy}
    \If{$pruned \neq \textit{None}$}
      \State $result \gets \Call{Backtrack}{assignment, pruned, C}$
      \If{$result \neq \textit{None}$} \Return $result$
      \EndIf
    \EndIf
    \State remove $X$ from $assignment$ \Comment{$domains$ unchanged $=$ restored}
  \EndFor
  \State \Return \textit{None}
\EndFunction
\Statex
\Function{Backtracking-Search}{$variables, domains, C$}
  \State $domains \gets$ each domain filtered by its latest-slot limit
  \If{some domain is empty} \Return \textit{None}
  \EndIf
  \State \Return \Call{Backtrack}{$\{\}, domains, C$}
\EndFunction
\end{algorithmic}
\end{algorithm}

On the sample files, the latest-slot filter leaves D2 and D3 with domain $\{1,2,3\}$. They tie on
MRV and on degree (one precedence neighbour each), so file order picks D2, and LCV places it in T1.
The resulting schedule is D2$\to$T1, D1$\to$T2, D3$\to$T3, D4$\to$T4. The notebook prints each
schedule together with a check of every constraint (Figures~\ref{fig:run1-schedule}, \ref{fig:run2-schedule} and~\ref{fig:run3-schedule}).
